\begin{abstract}
Augmentative releases of the parasitoid \textit{Tamarixia radiata} are a main tool against \textit{Diaphorina citri}, the vector of huanglongbing, but release models usually treat the orchard as a single patch. We extended a four-compartment model of psyllid, flush, parasitoid and plant vigour with impulsive releases to networks of eight patches coupled by diffusion of psyllids and parasitoids, with optional transit mortality, and compared uniform releases with releases targeted at the outbreak patch. Without transit mortality the amount released dominated the outcome. Uniform releases reduced the network pest load by about 30\% and targeted releases by 2.5--4.5\%, efficiencies per parasitoid were similar, and topology changed the pest load by less than 2\%, mainly through the degree of the outbreak patch. The master stability function ranked the networks by spatial coherence but did not account for its level over an outbreak. Parasitoid transit mortality created an intermediate dispersal rate that maximised the efficiency of targeted releases. This optimum maximises the pest load avoided relative to no release, not suppression, and lies where transit losses over the whole network drive the resident parasitoid towards collapse. A linear invasion argument explains why it scales with the mean degree and the mean edge mortality.
\end{abstract}
